\documentclass{cumcmthesis}

\usepackage{float}
\usepackage{subcaption}
\usepackage{booktabs}
\usepackage{bm}
\usepackage{url}

\title{基于状态空间偏差识别与任务事件调度的多源机器人定位融合及任务优化模型}
\contestname{数学建模竞赛论文}
\school{新疆大学}
\teamname{队伍编号}
\members{王雯瑾 \quad 梁嘉亮 \quad 李睿林}
\advisor{指导教师姓名}
\date{2026 年 5 月 2 日}

\newcommand{\argmin}{\mathop{\mathrm{arg\,min}}}
\newcommand{\norm}[1]{\left\lVert #1 \right\rVert}
\newcommand{\vfigpath}{figures/paper/v4}

\begin{document}
\maketitle

\begin{abstract}
本文围绕两源机器人定位数据的时间对齐、偏差识别和固定轨迹任务安排建立统一模型。对附件1，在无噪声假设下以方式1为参考搜索方式2时间偏差，得到 $\Delta=-198.4317\ \mathrm{s}$，重叠区间为 $270.4000$--$970.7500\ \mathrm{s}$，对齐后 RMSE 为 $9.06\times10^{-8}\ \mathrm{m}$，并输出 $7004$ 个 10 Hz 融合轨迹点。

对附件2，鲁棒中位数配准得到 $\Delta=-50.4947\ \mathrm{s}$，方式2相对方式1的固定偏差为 $(3.480,-1.822)\ \mathrm{m}$，RMSE 由 $3.997\ \mathrm{m}$ 降至 $0.7315\ \mathrm{m}$。进一步建立含状态 $[x,y,v_x,v_y,b_x,b_y]^\top$ 的卡尔曼滤波与 RTS 平滑模型；三组 $Q/R$ 参数下 $\Delta$ 完全一致，$b_x,b_y$ 变化小于 $0.01\ \mathrm{m}$，说明卡尔曼滤波与 RTS 平滑的主要作用不是显著降低 RMSE，而是验证固定偏差稳定性并提供平滑 10 Hz 状态轨迹。

对附件3，本文比较 M0--M4 偏差结构模型并采用 5 个连续时间块交叉验证。结果显示固定平移模型 M1 相比 M0 的 CV-RMSE 改善仅约 $0.10\%$，Kernel Ridge 等非固定模型训练误差下降但 CV 误差未稳定优于 M1。因此本文不做强校正，而将局部残差结构作为问题四任务鲁棒性风险来源。

对问题四，本文构造射击和拍照事件并建立联合 MILP。SciPy HiGHS 求得全局最优，联合方案选择 25 个事件，其中 21 次射击、4 次拍照，覆盖 9 个射击目标和 6 个拍照目标，射击期望命中数为 $8.0998$，综合目标效用为 $14.0998$。同时给出鲁棒裕度优先的保守方案：21 个事件、期望命中 $6.3993$、拍照覆盖 5 个目标、无低裕度风险事件。

\keywords{多源定位融合；卡尔曼滤波；RTS 平滑；系统偏差结构识别；混合整数规划；任务事件调度}
\end{abstract}

\section{问题重述与建模思路}
题目要求根据两种定位方式的轨迹数据完成时间对齐、10 Hz 融合输出、含噪数据偏差识别，以及在固定机器人轨迹下安排射击和拍照任务。定位部分的关键在于两源采样频率不同、时间轴偏移且误差结构不同；任务部分的关键在于射击和拍照共享时间轴，并同时受距离、速度、加速度和准备窗口约束。

本文采用“先轨迹融合、后事件调度”的路线。问题一解决无噪声双源轨迹对齐；问题二用鲁棒配准和状态空间模型识别固定偏差；问题三用时间块交叉验证判断偏差结构能否泛化；问题四将任务离散为可行事件，并用 MILP 同时处理射击概率收益、拍照覆盖和时间冲突。

\section{符号说明与基本假设}
记两源插值轨迹为 $\bm p_1(t),\bm p_2(t)$，方式2时间修正为 $t_{2,\mathrm{aligned}}=t_2+\Delta$。方式2相对方式1的固定偏差为 $\bm b=(b_x,b_y)^\top$。融合轨迹为 $\bm p_f(t)$，速度和加速度模长分别为 $v(t)$、$a(t)$。射击目标集合为 $\mathcal S$，拍照目标集合为 $\mathcal P$。

假设两种定位方式观测同一机器人；附件1无噪声，故不估计固定空间偏差；附件2的主要系统差异可由固定平移描述；附件3不预设偏差结构，必须用泛化误差判断；问题四中机器人轨迹固定，不进行路径重规划。

\section{问题一：无噪声双源轨迹的时间对齐与 10Hz 融合输出}
\subsection{两源插值与时间偏差搜索}
对第 $i$ 种定位方式构造三次样条插值
\begin{equation}
  \bm p_i(t)=\left(S_{i,x}(t),S_{i,y}(t)\right)^\top,\quad i=1,2.
\end{equation}
给定候选 $\Delta$ 后，方式2在方式1时间轴上的取值为 $\bm p_2(t-\Delta)$。在重叠区间 $\Omega(\Delta)$ 的 10 Hz 网格 $\mathcal T_\Delta$ 上定义
\begin{equation}
  J(\Delta)=\sqrt{\frac{1}{|\mathcal T_\Delta|}
  \sum_{t\in\mathcal T_\Delta}\norm{\bm p_2(t-\Delta)-\bm p_1(t)}_2^2}.
\end{equation}
附件1不估计空间偏差，直接以 $J(\Delta)$ 最小作为时间对齐准则。

\begin{table}[H]
  \centering
  \caption{附件1时间对齐与 10 Hz 输出摘要}
  \label{tab:attach1}
  \small
  \begin{tabular}{lrrrrrr}
    \toprule
    数据 & $\Delta$/s & 重叠起点/s & 重叠终点/s & 10Hz点数 & 对齐前RMSE/m & 对齐后RMSE/m\\
    \midrule
    附件1 & -198.4317 & 270.4000 & 970.7500 & 7004 & $9.06\times10^{-8}$ & $9.06\times10^{-8}$\\
    \bottomrule
  \end{tabular}
\end{table}
10 Hz 轨迹输出为 \texttt{outputs/trajectories/fused\_attachment1\_10hz.csv}，独立汇总表为 \texttt{outputs/tables/attachment1\_alignment\_summary.csv}。

\begin{figure}[H]
  \centering
  \includegraphics[width=0.96\textwidth]{\vfigpath/attachment1_time_alignment.pdf}
  \caption{附件1时间对齐、RMSE 搜索曲线与 10 Hz 融合轨迹}
  \label{fig:attach1}
\end{figure}

\section{问题二：含噪数据的去噪、固定偏差估计与状态空间融合}
\subsection{鲁棒配准基准}
附件2存在随机噪声和固定偏差。对每个 $\Delta$，先计算残差 $\bm r(t)=\bm p_2(t-\Delta)-\bm p_1(t)$，再用坐标中位数估计固定偏差 $\bm b(\Delta)$。目标函数为校正后的 RMSE。该方法对局部离群噪声不敏感，并作为状态空间模型的初始化基准。

\subsection{卡尔曼滤波与 RTS 平滑}
首次使用时称为 Kalman Filter（卡尔曼滤波）和 Rauch--Tung--Striebel smoother（RTS 平滑），后文统一用中文。状态向量取
\begin{equation}
  \bm X=[x,\ y,\ v_x,\ v_y,\ b_x,\ b_y]^\top .
\end{equation}
其中 $b_x,b_y$ 表示方式2相对方式1的固定偏差；随机观测误差不是状态变量，而由观测噪声协方差 $R_1,R_2$ 描述。状态转移采用常速度模型，偏差过程噪声设为很小值。观测矩阵为
\begin{equation}
  \bm H_1=
  \begin{bmatrix}1&0&0&0&0&0\\0&1&0&0&0&0\end{bmatrix},\quad
  \bm H_2=
  \begin{bmatrix}1&0&0&0&1&0\\0&1&0&0&0&1\end{bmatrix}.
\end{equation}
方式1作为参考，不估计其固定偏差；方式2观测中显式包含 $\bm b$。

\subsection{参数灵敏度与结果分析}
\begin{table}[H]
  \centering
  \caption{附件2鲁棒基准、卡尔曼滤波与 RTS 平滑结果}
  \label{tab:kalman_compare}
  \small
  \begin{tabular}{lrrrrr}
    \toprule
    方法 & $\Delta$/s & $b_x$/m & $b_y$/m & 对齐前RMSE/m & 校正后RMSE/m\\
    \midrule
    RobustMedian & -50.4947 & 3.4800 & -1.8221 & 3.9970 & 0.7315\\
    KalmanFilter & -50.4947 & 3.4809 & -1.8284 & 3.9970 & 0.7314\\
    KalmanRTS & -50.4947 & 3.4754 & -1.8337 & 3.9970 & 0.7314\\
    \bottomrule
  \end{tabular}
\end{table}

\begin{table}[H]
  \centering
  \caption{附件2 Kalman 模型 $Q/R$ 参数灵敏度}
  \label{tab:qr}
  \scriptsize
  \begin{tabular}{lrrrrrrrr}
    \toprule
    配置 & $q_p$ & $q_v$ & $q_b$ & $r_{1x}$ & $r_{2x}$ & $\Delta$/s & $b_x$/m & $b_y$/m\\
    \midrule
    low\_process & 0.002 & 0.010 & $10^{-7}$ & 0.2675 & 0.2675 & -50.4947 & 3.4807 & -1.8283\\
    balanced & 0.010 & 0.050 & $10^{-6}$ & 0.2675 & 0.2675 & -50.4947 & 3.4803 & -1.8275\\
    smooth\_bias & 0.030 & 0.100 & $10^{-8}$ & 0.3210 & 0.3210 & -50.4947 & 3.4809 & -1.8284\\
    \bottomrule
  \end{tabular}
\end{table}
三组配置下 $\Delta$ 完全一致，固定偏差变化小于 $0.01\ \mathrm{m}$，RMSE 变化不足 $10^{-4}\ \mathrm{m}$。因此卡尔曼滤波与 RTS 平滑的主要贡献是从状态空间角度验证固定偏差估计稳定性，并输出后续速度、加速度计算所需的平滑状态轨迹，而不是包装为显著降低 RMSE 的性能模型。

\begin{figure}[H]
  \centering
  \begin{subfigure}{0.32\textwidth}
    \includegraphics[width=\textwidth]{\vfigpath/attachment2_method_rmse_compare.pdf}
    \caption{RMSE 对比}
  \end{subfigure}
  \begin{subfigure}{0.32\textwidth}
    \includegraphics[width=\textwidth]{\vfigpath/attachment2_trajectory_kalman_rts.pdf}
    \caption{轨迹对比}
  \end{subfigure}
  \begin{subfigure}{0.32\textwidth}
    \includegraphics[width=\textwidth]{\vfigpath/attachment2_residual_distribution_compare.pdf}
    \caption{残差分布}
  \end{subfigure}
  \caption{附件2鲁棒配准与 RTS 平滑验证型增强结果}
  \label{fig:attach2}
\end{figure}

\section{问题三：实际测量数据的系统偏差结构识别}
\subsection{残差特征构造}
附件3先做时间对齐，得到 $\Delta=367.8843\ \mathrm{s}$。原始插值残差 RMSE 约为 $2.750\ \mathrm{m}$；表~\ref{tab:bias_structure} 中的 RMSE 基于平滑状态残差计算，约为 $1.94\ \mathrm{m}$，两者口径不同：前者评价原始数据对齐效果，后者用于偏差结构比较。

残差特征包括归一化时间、融合位置、速度、加速度、航向角正余弦和曲率近似：
\begin{equation}
  \phi=(t,x,y,v,a,\cos\theta,\sin\theta,\kappa).
\end{equation}

\subsection{M0--M4 偏差结构模型}
\begin{align}
  \mathrm{M0}:&\quad \hat{\bm r}=0,\\
  \mathrm{M1}:&\quad \hat{\bm r}=\bm b_0,\\
  \mathrm{M2}:&\quad \hat r_x=\alpha_{0x}+\alpha_{1x}t,\quad \hat r_y=\alpha_{0y}+\alpha_{1y}t,\\
  \mathrm{M3}:&\quad \hat{\bm r}=W\phi(t,x,y,v,a,\cos\theta,\sin\theta,\kappa)+\bm c,\\
  \mathrm{M4}:&\quad \hat{\bm r}=f_{\mathrm{KRR}}(\phi).
\end{align}
其中 M3 用 Ridge 线性特征映射，M4 用 RBF Kernel Ridge，模型复杂度按训练样本规模记录。

\subsection{时间块交叉验证}
将时间序列按顺序分为 $K=5$ 个连续块。每次取 $K-1$ 块训练、1 块测试；标准化参数只在训练块计算并应用到测试块；Ridge 和 Kernel Ridge 的超参数在建模前固定，不使用测试块调参。

\begin{table}[H]
  \centering
  \caption{附件3系统偏差结构模型比较}
  \label{tab:bias_structure}
  \small
  \begin{tabular}{llrrrr}
    \toprule
    模型 & 含义 & 复杂度 & 训练RMSE/m & 时间块CV-RMSE/m & 相对M1改善\\
    \midrule
    M0 & 无系统偏差 & 0 & 1.9456 & 1.9456 & 0.0000\\
    M1 & 固定平移偏差 & 2 & 1.9402 & 1.9436 & 0.0000\\
    M2 & 线性时变偏差 & 4 & 1.9379 & 1.9463 & -0.0014\\
    M3 & 状态相关线性偏差 & 16 & 1.9159 & 2.1733 & -0.1182\\
    M4 & Kernel Ridge 非线性偏差 & 1500 & 1.6217 & 2.0319 & -0.0454\\
    \bottomrule
  \end{tabular}
\end{table}

\subsection{结论与融合策略}
附件3残差存在局部结构性波动，但固定平移和非固定偏差模型在时间块交叉验证下均未形成稳定可泛化的强校正模型。全局固定平移偏差在工程幅值和时间块 Bootstrap 下也不构成稳定强校正依据。因此本文不采用 M2--M4 做强校正，而将残差结构作为任务鲁棒性风险来源，在问题四中通过裕度与保守方案吸收该风险。

\begin{figure}[H]
  \centering
  \begin{subfigure}{0.48\textwidth}
    \includegraphics[width=\textwidth]{\vfigpath/attachment3_residual_time_trend_clean.pdf}
    \caption{残差时间趋势}
  \end{subfigure}
  \begin{subfigure}{0.48\textwidth}
    \includegraphics[width=\textwidth]{\vfigpath/attachment3_residual_vector_field_clean.pdf}
    \caption{残差向量场}
  \end{subfigure}
  \caption{附件3残差局部结构诊断}
\end{figure}

\begin{figure}[H]
  \centering
  \begin{subfigure}{0.48\textwidth}
    \includegraphics[width=\textwidth]{\vfigpath/attachment3_bias_model_cv_compare_clean.pdf}
    \caption{训练与时间块 CV}
  \end{subfigure}
  \begin{subfigure}{0.48\textwidth}
    \includegraphics[width=\textwidth]{\vfigpath/attachment3_residual_state_relation_clean.pdf}
    \caption{残差与状态关系}
  \end{subfigure}
  \caption{附件3偏差结构模型与状态特征证据}
\end{figure}

\section{问题四：射击与拍照任务事件优化}
\subsection{射击事件与命中概率}
射击事件 $e=(k,t)$ 在准备窗口 $[t-1.5,t]$ 内需满足
\begin{equation}
  5\le d_k(\tau)\le 30,\quad v(\tau)\le 2,\quad a(\tau)\le 1.5.
\end{equation}
设二元变量 $x_e$ 表示是否选择射击事件，$n_k=\sum_{e\in\mathcal E_k^S}x_e$。第 $m$ 次射击的边际收益为
\begin{equation}
  g_m=0.85\cdot 0.15^{m-1},
\end{equation}
目标 $k$ 至少命中一次的概率为
\begin{equation}
  P_k(n_k)=1-0.15^{n_k}.
\end{equation}
用二元变量 $y_{k,m}$ 线性化递减收益，并加约束 $y_{k,m}\le y_{k,m-1}$、$\sum_m y_{k,m}=n_k$。

\subsection{拍照事件与多目标覆盖}
拍照事件 $e=(t,\phi)$ 的准备窗口为 $[t-0.5,t]$，需满足
\begin{equation}
  10\le d_p(\tau)\le 40,\quad v(\tau)\le 1.5,\quad a(\tau)\le 1.5.
\end{equation}
若目标 $p$ 落入视场角内，则 $p\in C_e$。二元变量 $z_e$ 表示拍照事件是否选择，$u_p$ 表示拍照目标是否至少被覆盖一次，并加
\begin{equation}
  u_p\le \sum_{e:p\in C_e}z_e .
\end{equation}
若两个拍照事件都覆盖同一目标且观测方向圆周差小于 $60^\circ$，则加约束 $z_i+z_j\le1$，其中角差为
\begin{equation}
  \min(|\theta_i-\theta_j|,\ 360^\circ-|\theta_i-\theta_j|).
\end{equation}

\subsection{联合 MILP 调度模型}
所有事件统一用变量 $q_e$ 表示。若事件准备窗口重叠，则
\begin{equation}
  q_i+q_j\le1.
\end{equation}
目标函数为
\begin{equation}
 \max\ M_1(E_{\mathrm{shoot}}+C_{\mathrm{photo}})
 +M_2O_{\mathrm{photo}}+M_3B_{\mathrm{multi}}
 +M_4\sum_e m_eq_e-M_5R_{\mathrm{risk}},
\end{equation}
其中 $E_{\mathrm{shoot}}=\sum_{k,m}g_my_{k,m}$，$C_{\mathrm{photo}}=\sum_pu_p$，$O_{\mathrm{photo}}=\sum_e |C_e|z_e$，$B_{\mathrm{multi}}=\sum_e\max(0,|C_e|-1)z_e$。本文采用 $M_1\gg M_2\gg M_3\gg M_4\gg M_5$，但将射击命中和拍照覆盖放在同一级目标完成收益中，避免射击单独压倒拍照。

\begin{table}[H]
  \centering
  \caption{联合 MILP 求解状态}
  \label{tab:solver}
  \scriptsize
  \begin{tabular}{llrrrrr}
    \toprule
    方案 & 求解器 & 状态 & 目标值 & gap & 变量数 & 约束数\\
    \midrule
    收益优先 & scipy.milp/HiGHS & Optimal & 14163.1600 & 0.0 & 1111 & 44294\\
    裕度优先 & scipy.milp/HiGHS & Optimal & 11451.9716 & 0.0 & 800 & 22583\\
    \bottomrule
  \end{tabular}
\end{table}

\subsection{纯射击、纯拍照与联合方案比较}
\begin{table}[H]
  \centering
  \caption{纯射击、纯拍照与联合方案比较}
  \label{tab:plan_compare}
  \small
  \begin{tabular}{lrrrrrrrr}
    \toprule
    方案 & 射击事件 & 拍照事件 & 总事件 & 射击覆盖 & 拍照覆盖 & 期望命中 & 拍照观测 & 综合效用\\
    \midrule
    纯射击 & 21 & 0 & 21 & 9 & 0 & 8.0998 & 0 & 8.0998\\
    纯拍照 & 0 & 3 & 3 & 0 & 5 & 0.0000 & 5 & 5.0000\\
    联合MILP & 21 & 4 & 25 & 9 & 6 & 8.0998 & 6 & 14.0998\\
    \bottomrule
  \end{tabular}
\end{table}
联合方案相比纯射击几乎不损失射击期望收益，同时额外获得 6 个拍照目标覆盖。

\begin{table}[H]
  \centering
  \caption{收益优先方案与保守方案比较}
  \label{tab:conservative}
  \small
  \begin{tabular}{lrrrrrr}
    \toprule
    方案 & 总事件 & 期望命中 & 拍照覆盖 & 综合效用 & 最小裕度 & 风险数\\
    \midrule
    A 收益优先 & 25 & 8.0998 & 6 & 14.0998 & 0.0031 & 5\\
    B 裕度优先 & 21 & 6.3993 & 5 & 11.3993 & 0.0337 & 0\\
    \bottomrule
  \end{tabular}
\end{table}
A 用于最大收益；B 通过剔除低于 $0.03$ 的边界事件并提高风险惩罚，适合工程保守场景。

\subsection{逐目标结果与最终任务表}
完整逐目标表已输出到 \texttt{outputs/tables/shooting\_target\_summary.csv} 和 \texttt{outputs/tables/photo\_target\_summary.csv}。射击方面，S03、S06、S10--S16 等 9 个目标被选择，不可达目标如 S07--S09、S17、S18 标记为 \texttt{no\_feasible\_window}；未选但有窗口的目标主要受时间冲突影响。拍照方面，P01、P02、P03、P04、P10、P11 被覆盖，其余多为 \texttt{no\_feasible\_window}。

\begin{table}[H]
  \centering
  \caption{联合 MILP 最终任务事件表}
  \label{tab:final_events}
  \scriptsize
  \begin{tabular}{rlllrrr}
    \toprule
    序号 & 类型 & 目标/覆盖目标 & 事件ID & 开始/s & 执行/s & 裕度\\
    \midrule
    1 & 射击 & S06 & S\_00169 & 479.45 & 480.95 & 0.034\\
    2 & 射击 & S03 & S\_00106 & 481.55 & 483.05 & 0.116\\
    3 & 拍照 & P11 & P\_00001 & 493.05 & 493.55 & 0.011\\
    4 & 拍照 & P03,P04 & P\_00012 & 494.05 & 494.55 & 0.217\\
    5 & 拍照 & P10 & P\_00020 & 495.05 & 495.55 & 0.072\\
    6 & 射击 & S12 & S\_00665 & 506.75 & 508.25 & 0.009\\
    7 & 射击 & S11 & S\_00457 & 508.55 & 510.05 & 0.314\\
    8 & 拍照 & P01,P02 & P\_00035 & 510.55 & 511.05 & 0.184\\
    9 & 射击 & S12 & S\_00719 & 513.75 & 515.25 & 0.087\\
    10 & 射击 & S11 & S\_00509 & 515.35 & 516.85 & 0.278\\
    11 & 射击 & S12 & S\_00751 & 516.95 & 518.45 & 0.192\\
    12 & 射击 & S11 & S\_00526 & 522.15 & 523.65 & 0.125\\
    13 & 射击 & S11 & S\_00542 & 523.75 & 525.25 & 0.108\\
    14 & 射击 & S10 & S\_00306 & 525.35 & 526.85 & 0.147\\
    15 & 射击 & S10 & S\_00322 & 526.95 & 528.45 & 0.085\\
    16 & 射击 & S10 & S\_00333 & 531.45 & 532.95 & 0.121\\
    17 & 射击 & S10 & S\_00350 & 533.15 & 534.65 & 0.062\\
    18 & 射击 & S10 & S\_00366 & 534.75 & 536.25 & 0.108\\
    19 & 射击 & S15 & S\_00910 & 748.85 & 750.35 & 0.056\\
    20 & 射击 & S14 & S\_00889 & 750.45 & 751.95 & 0.029\\
    21 & 射击 & S16 & S\_00979 & 752.05 & 753.55 & 0.025\\
    22 & 射击 & S13 & S\_00866 & 762.05 & 763.55 & 0.003\\
    23 & 射击 & S12 & S\_00827 & 763.65 & 765.15 & 0.100\\
    24 & 射击 & S11 & S\_00644 & 765.55 & 767.05 & 0.080\\
    25 & 射击 & S12 & S\_00865 & 767.45 & 768.95 & 0.107\\
    \bottomrule
  \end{tabular}
\end{table}
填写模板输出为 \texttt{outputs/result\_filled\_v3.xlsx}，未覆盖旧版 \texttt{v2} 文件。

\begin{figure}[H]
  \centering
  \begin{subfigure}{0.48\textwidth}
    \includegraphics[width=\textwidth]{\vfigpath/joint_task_gantt.pdf}
    \caption{事件 Gantt 图}
  \end{subfigure}
  \begin{subfigure}{0.48\textwidth}
    \includegraphics[width=\textwidth]{\vfigpath/joint_task_spatial_distribution.pdf}
    \caption{空间分布}
  \end{subfigure}
  \caption{联合任务方案的时间与空间分布}
\end{figure}

\section{模型评价与灵敏度分析}
问题一在无噪声附件上实现了几乎零 RMSE 的时间对齐，并明确给出 10 Hz 输出。问题二中，RTS 平滑与鲁棒基准数值一致，说明固定偏差估计稳定；其价值在于验证和状态平滑。问题三避免将局部残差波动过度解释为可泛化强校正模型。问题四用 MILP 坐实“优化”，并给出收益优先和保守裕度优先两种可答辩方案。

\section{结论}
本文补齐了附件1无噪声时间对齐与 10 Hz 输出，修正了附件2卡尔曼滤波与 RTS 平滑的模型定位，完成了附件3偏差结构识别的时间块 CV 论证，并将问题四从贪心启发式升级为事件级 MILP。最终联合方案在保持纯射击期望命中 $8.0998$ 的同时覆盖 6 个拍照目标；保守方案牺牲部分收益换取更高裕度，可作为工程备选。

\begin{thebibliography}{12}
\bibitem{kalman} Kalman R E. A New Approach to Linear Filtering and Prediction Problems[J]. Journal of Basic Engineering, 1960, 82(1): 35--45.
\bibitem{rts} Rauch H E, Tung F, Striebel C T. Maximum likelihood estimates of linear dynamic systems[J]. AIAA Journal, 1965, 3(8): 1445--1450.
\bibitem{efron} Efron B, Tibshirani R J. An Introduction to the Bootstrap[M]. New York: Chapman and Hall/CRC, 1993.
\bibitem{hastie} Hastie T, Tibshirani R, Friedman J. The Elements of Statistical Learning[M]. New York: Springer, 2009.
\bibitem{wolsey} Wolsey L A. Integer Programming[M]. New York: Wiley, 1998.
\bibitem{nemhauser} Nemhauser G L, Wolsey L A. Integer and Combinatorial Optimization[M]. New York: Wiley, 1988.
\bibitem{bertsimas} Bertsimas D, Sim M. The Price of Robustness[J]. Operations Research, 2004, 52(1): 35--53.
\bibitem{boyd} Boyd S, Vandenberghe L. Convex Optimization[M]. Cambridge: Cambridge University Press, 2004.
\end{thebibliography}

\end{document}
